
\def \Subject {تقسیم و حل}
\def \Author {زهرا محقق}
\def \Date {۱۳۹۸/۷/۱۸}
\def \Session {5}
\setcounter{chapter}{\Session-1}

\chapter{\Subject}
\chapterauthor{\Author~ - \Date}

جزوه جلسه
\Session ام
 مورخ 
\Date
درس 
\CourseName
 تهیه شده توسط
 \Author.
در جهت مستند کردن مطالب درس 
\CourseName
، بر آن شدیم که از دانشجویان جهت مکتوب کردن مطالب کمک بگیریم. هر دانشجو می‌تواند برای مکتوب کردن یک جلسه داوطلب شده و با توجه به کیفیت جزوه از لحاظ کامل بودن مطالب، کیفیت نوشتار و استفاده از اشکال و منابع کمک آموزشی، حداکثر یک نمره مثبت از بیست نمره دریافت کند. خواهش‌مند است نام و نام خانوادگی خود، عنوان درس، شماره و تاریخ جلسه در ابتدای این فایل را با دقت پر کنید. مطالبی که در ادامه آمده فقط جنبه راهنمایی شیوه استفاده از لاتک می‌باشد. خواهشمند است این پاراگراف 

همچنین برای اضافه کردن شکل می‌توانید از شکل زیر استفاده کنید و برای ارجاع دادن به بصورت شکل

 استفاده کنید.
 همچنین برای درج کلمات انگلیسی در پاراگراف فارسی می‌توان به این شکل

 یا برای تاکید به این شکل
 
همانطور که در نمونه‌کدها ملاحظه می‌کنید فاصله‌ها دقیقا همانطور که در لاتک آمده است در خروجی نمایش داده می‌شود.
 


لذا لازم است که فاصله‌های ابتدای خط را مرتب و منظم و بدون کوچکترین اشکالی رعایت کنید تا خروجی مستند شما نیز بدون اشکال باشد.


در نهایت استفاده از دستور 
\grayBox{grayBox}
نیز می‌تواند کمک شایانی به زیبایی و خوانایی متن شما بکند. این دستور علاوه بر عوض کردن رنگ پس‌زمینه از فونت انگلیسی با عرض ثابت نیز استفاده می‌کند که برای کلمات کلیدی یا اسامی متغیرها در کد قابل استفاده است.



همچنین سعی کنید حتی‌الامکان به منابع و مراحع مناسب ارجاع دهید
\cite{CLRS}. 
علاوه بر مراجع چنانچه ابزار یا وب‌سایت قابل توجهی موجود است خوب است به آن هم ارجاع دهید
\cite{visualizationwebsite}. 

